% =====================================================================
% 0. AI assistant environment (initial section required by the course)
% =====================================================================
\section{AI assistant environment}\label{sec:ai-env}

The course asks MLLB reports to open with the AI environment used and how it was controlled. The
record of every interaction is \texttt{docs/AI\_USAGE\_LOG.md} in the repository; the errors found are
in Appendix~\ref{app:ai-code}.

\subsection{Tools and roles}

\begin{table}[htbp]
\centering
\caption{AI assistants and supporting tools used in the project.}
\label{tab:ai-tools}
\small
\begin{tabularx}{\textwidth}{>{\raggedright\arraybackslash}p{3.4cm} >{\raggedright\arraybackslash}X >{\raggedright\arraybackslash}X}
\toprule
\textbf{Tool} & \textbf{Role in the project} & \textbf{Main limitation observed} \\
\midrule
Claude (Anthropic), desktop app with access to the project folder & Proposals of every design decision
with its reasons; code of \texttt{src/} and \texttt{tests/}; figures and tables generated by code; block
notes; drafts of the reports. Tests run on a CPU copy with synthetic data before delivery &
Errors visible only on the real GPU or the real data (out of memory, tensors on two devices); files
reported as written when an older copy had reached the folder; one claim stated from memory
(Appendix~\ref{app:ai-code}) \\
Perplexity & Candidate papers per topic, without quotes or numbers & --- (used only to find PDFs) \\
NotebookLM, restricted to the downloaded PDFs & Literal quotes with file and page for each claim &
Wrong pages, mixed versions of a paper, quotes that do not support the claim \\
\midrule
VS Code, Python 3.12 virtual environment, PyTorch with CUDA 12.6 & Execution of every run, test and
report step by us, on Oriana's laptop (NVIDIA RTX 3050 Ti, 4\,GB) & --- \\
git and GitHub & Version control; commits and pushes made by us & --- \\
Overleaf & Compilation of the reports & --- \\
\bottomrule
\end{tabularx}
\end{table}

\subsection{How the assistant was controlled}

The main risk of an assistant that writes complete pipelines is accepting results that one cannot
explain. Seven rules limited it:
\begin{enumerate}[leftmargin=1.4em,itemsep=1pt]
  \item \textbf{Decisions before code.} Every block started with a written proposal of the decisions
        (architecture, ablation grid, seeds, selection rule, outputs) with its reasons and course
        sources; we accepted, changed or rejected it, and only then was code written. The accepted
        version is a dated entry of \texttt{docs/DECISIONS.md}; a decision changes only with a new dated
        entry.
  \item \textbf{We execute.} The assistant never trained on the real data. It tested its code on a CPU
        copy with synthetic data; we ran the tests, the smoke runs and the real runs on the laptop and
        reported the output. No result counts until we have run it.
  \item \textbf{No number typed by hand.} Every number of the reports is a macro generated from the run
        logs (Section~\ref{sec:code}); the only numbers typed by hand are the published benchmark values,
        each stored with its file, page and quote.
  \item \textbf{No claim without a checked quote.} Literature claims come only from quotes checked
        against the PDFs; interpretations without a number behind them are written as hypotheses.
  \item \textbf{Every figure inspected} against the numbers it shows before it was accepted.
  \item \textbf{Every write checked.} After several cases in which an older copy of a file reached the
        folder, every file written by the assistant is compared with the intended version (size or
        content) before it is reported as written.
  \item \textbf{One block per conversation.} The work was split into one conversation per block
        (data, EDA, evaluation, one per architecture, learning strategies, literature, reports), which
        share nothing but the project folder. Each ends with a block note in \texttt{docs/blocks/},
        updated plan and log, and our commit.
\end{enumerate}
In my own blocks I re-read every line of code with the assistant's explanation of each technique before
running it, because the oral defence asks about the theory and about specific lines.

\subsection{Prompt strategy and multi-step interaction}

Each conversation of the project opened with a short, fixed prompt that gave the assistant its scope and
forced a plan before any code. The start prompt of my generative block, translated from Spanish:
\begin{quote}\small\itshape
You are the chat ``DL 13 $\cdot$ Generative (ECG-GPT)'' of the Deep Learning project (PTB-XL) of Oriana
Ferraz and Mar\'ia Mart\'inez. Working folder: C:\textbackslash Users\textbackslash ori\textbackslash DL.
Before anything else read CHATS\_DL.md (common protocol and card 13), DECISIONS.md, PLAN\_DL.md and
HANDOFF.md. Then tell us in a few lines what you have understood, what is missing and where you start.
Do not write code or documents until we confirm.
\end{quote}
From there a block went through the same steps: the assistant proposed the decisions with their course
sources; we accepted or changed them and they were written to \texttt{docs/DECISIONS.md}; the assistant
wrote the code and its tests; I ran the tests, a smoke run and the real runs and pasted the output back;
the assistant analysed the logs and wrote the block note; and the conversation closed with the plan, the
log of \texttt{docs/AI\_USAGE\_LOG.md} and our commit. When an output was wrong, the next prompt contained
the error message or the file listing, never a guess (Appendix~\ref{app:ai-code}).
